% Folder 91, batch 08 — archivists' pages 141–142.
% Pass: Opus 5.5 (claude-opus-5-5), 2026-10-03 — first pass, unchecked against the pages by a human.
%
% Page 142 is an unrelated typed verso (a typescript of Lie-group exercises,
% scanned upside down) carrying nothing in his hand: skipped, per the rule.
\documentclass[11pt,a4paper]{article}
\input{../preamble/grothendieck}

\folder{91}
\batch{8}
\pages{141}{142}
\foldertitle{Autour de Néron / Greenberg-Néron. Foncteurs Hom (méthodes non-projectives) : notes manuscrites (s.d.), lettre (1967).}
\dating{[à partir de 1964-vers 1970]}
\watermark{Édition de démonstration}

\begin{document}

\page{141}
\note{la page s'ouvre sur une liste dont le début n'est pas dans ce lot ; l'argument vient vraisemblablement des pages précédentes du dossier.}

\struck{\uncertain{$\mathfrak{m}/\mathfrak{m}^{2}$}}

\ill{} \ill{} \ill{} \uncertain{$\mathcal{T}^{0}_{B/A_{0}} \otimes V_{1}$}

\ill{} \ill{} \ill{} \uncertain{$\mathcal{T}^{0}_{B_{0}/A_{0}} \otimes V_{2}$}

\note{à la ligne suivante, un trait horizontal tient lieu des mots, comme un « idem ».}

\uncertain{$\mathcal{T}^{0}_{B_{0}/A_{0}} \otimes V_{3}$} \qquad \uncertain{$\mathcal{T}_{V_{0}}/$}

\note{la lettre lue $\mathcal{T}$ est un V cursif dont le jambage droit descend en crochet ; la lecture comme faisceau tangent est une conjecture, suggérée par les produits tensoriels $\otimes V_{i}$ et les exposants $0$. Les trois premiers mots de chacune des deux premières lignes, très abrégés, n'ont pu être lus.}

$\mathbb{G}_{A,n}$ = \uncertain{groupe pro-algébrique} sur $k$ des \uncertain{changements de base} de \ill{} \ill{}
\[
\begin{cases}
  f_{i} = a_{i} + \sum_{j} b_{ij}\, t_{j} + \sum_{2 \leq |p| \leq N} c_{i,p}\, t^{p},
  \qquad a_{i},\, b_{ij},\, c_{i,p} \in A, \\
  a_{i} \in \mathfrak{m}, \quad \det(b_{ij}) \notin \mathfrak{m}.
\end{cases}
\]

$\mathbb{G}_{n}$ \uncertain{groupe} sur $\mathbb{Z}$ des \uncertain{changements} de \ill{}
\[
\begin{cases}
  f_{i} = a_{i} + \sum b_{ij}\, t_{j} + \sum_{2 \leq |p| \leq N} c_{i,p}\, t_{1}^{p_{1}} \cdots t_{n}^{p_{n}}, \\
  \det(b_{ij}) \neq 0
\end{cases}
\]
i.e. des \uncertain{automorphismes} de \ill{}

\struck{\uncertain{$a + t$}} \quad $a_{0} + a_{1} t + a_{2} t^{2}$

\struck{$a'_{0} + a'_{1}(a_{0} + a_{1} t)$ \ill{}} $=$
\[
a'_{0} + a'_{1}(a_{0} + a_{1} t + a_{2} t^{2}) + a'_{2}(a_{0} + a_{1} t + a_{2} t^{2})^{2}
\]
\note{dans le carré, il souligne $a_{0} + a_{1} t$ et surligne $a_{1} t + a_{2} t^{2}$, repères pour le développement.}

$k[t]/(t^{3})$
\[
\begin{aligned}
  &= (a'_{0} + a'_{1} a_{0} + a'_{2} a_{0}^{2}) \\
  &\quad + (a'_{1} a_{1} + 2 a'_{2} a_{0} a_{1})\, t \\
  &\quad + (a'_{1} a_{2} + a'_{2} a_{1}^{2} + 2 a'_{2} a_{0} a_{2})\, t^{2} \\
  &\quad + 2 a'_{2} a_{1} a_{2}\, t^{3} \\
  &\quad + a'_{2} a_{2}^{2}\, t^{4}
\end{aligned}
\]
\note{les trois premières lignes (jusqu'au terme en $t^{2}$) sont encadrées, avec en regard « $k[t]/(t^{3})$ » : c'est la composition des deux polynômes modulo $t^{3}$. Le facteur $2$ de $2 a'_{2} a_{0} a_{1}$ est ajouté par lui au-dessus du $+$ ; dans $2 a'_{2} a_{0} a_{2}$ l'indice du $a'$ est mal formé et pourrait se lire $1$, la valeur $2$ est celle que demande le calcul. La page s'arrête sur ce développement ; la suite, s'il y en a une, n'est pas dans ce lot.}

\end{document}
